\documentclass[a4paper,10pt]{report}\input{../0.Préambule/Préambule_cours_prof}\begin{document}

\pagestyle{empty}
\newpage
\noindent NOM, Prénom et classe : \dotfill

\noindent \begin{minipage}{11cm}
\section*{Puissances - Notation scientifique}
\addcontentsline{toc}{section}{Évaluation : les puissances et la notation scientifique}

\noindent \textit{Il sera tenu compte dans la notation de la propreté ainsi que de la justification apportée à chacune des réponses.}

\noindent \textit{Le barème est donné à titre indicatif. Il pourra être modifié ultérieurement.}

\noindent \textit{L'usage de la calculatrice est \textbf{interdite}.} \hfill \textit{Durée : 30 minutes}
\end{minipage}
\hspace{6mm}
\begin{minipage}{5cm}
\renewcommand{\arraystretch}{1.5}
\begin{tabular}{|l|r|}
\hline
\textbf{1. Chercher :} & $\hspace{4mm}/\hspace{4mm}$\\\hline
\textbf{2. Modéliser :} & $\hspace{4mm}/\hspace{4mm}$\\\hline
\textbf{3. Représenter :} & $\hspace{4mm}/\hspace{4mm}$\\\hline
\textbf{4. Calculer :} & $\hspace{4mm}/\hspace{4mm}$\\\hline
\textbf{5. Raisonner :} & $\hspace{4mm}/\hspace{4mm}$\\\hline
\textbf{6. Communiquer :}$\hspace{2mm}$ & $\hspace{4mm}/\hspace{4mm}$\\\hline
\end{tabular}
\end{minipage}

\medskip
\paragraph{Question de cours}: \hfill \textbf{2 points}

\noindent Calculer chaque nombre :
\begin{enumerate}
\begin{tabularx}{\linewidth}{*{2}{>{\arraybackslash}X}}
\item $2^3$ = \dots ~\dots ~\dots &
\item $3^0$ = \dots ~\dots ~\dots \\
\end{tabularx}
\end{enumerate}

\medskip\noindent\grandscarreaux{\linewidth}{2.4cm}

\medskip
\paragraph{Exercice 1}: \hfill \textbf{4 points}

\noindent Compléter chaque égalité à l'aide d'un exposant entier relatif.
\begin{enumerate}
\begin{tabularx}{\linewidth}{*{4}{>{\arraybackslash}X}}
\item $10000$ = $10^{\dots ~\dots}$ &
\item $0,01$ = $10^{\dots ~\dots}$ &
\item $-0,0001$ = $-10^{\dots ~\dots}$ &
\item $-1000$ = $-10^{\dots ~\dots}$ \\
\end{tabularx}
\end{enumerate}

\medskip\noindent\grandscarreaux{\linewidth}{4cm}

\medskip
\paragraph{Exercice 2}: \hfill \textbf{3 points}

\noindent Écrire chaque expression sous la forme $3^n$, où $n$ un entier relatif.

\begin{enumerate}
\begin{tabularx}{\linewidth}{*{3}{>{\arraybackslash}X}}
\item $3^2 \times 3^5$ &
\item $\frac{3^7}{3^{-5}}$ &
\item $\left(3^{-5}\right)^{-2}$ \\
\end{tabularx}
\end{enumerate}

\medskip\noindent\grandscarreaux{\linewidth}{4.8cm}

\newpage
\paragraph{Exercice 3}: \hfill \textbf{4 points}

\noindent Écrire chaque expression sous la forme $a^n$, où $a$ est un nombre relatif et $n$ un entier relatif.
\begin{enumerate}
\begin{tabularx}{\linewidth}{*{2}{>{\arraybackslash}X}}
\item $\dfrac{2^{-3}\times 2^4 \times 2^{-2}}{2^5 \times 2^8}$ &
\item $\dfrac{7 \times 7^{26}}{\left( 7^5 \right) ^6 \times 7^{-7}}$ \\
\end{tabularx}
\end{enumerate}

\medskip\noindent\grandscarreaux{\linewidth}{5.6cm}

\medskip
\paragraph{Exercice 4}: \hfill \textbf{3 points}

\noindent Donner l'écriture scientifique de chaque nombre.
\begin{enumerate}
\begin{tabularx}{\linewidth}{*{3}{>{\arraybackslash}X}}
\item $\np{348000000}$ &
\item $\np{-0,547}$ &
\item $\np{0,0024}$ \\
\end{tabularx}
\end{enumerate}

\medskip\noindent\grandscarreaux{\linewidth}{5.6cm}

\medskip
\paragraph{Exercice 5}: \hfill \textbf{4 points}

\noindent Une analyse chimique a montré qu'il y avait $120$mg de magnésium dans $5$L d'eau.

\medskip
\noindent Calcule la concentration, en g/L, de magnésium dans cette eau, en expliquant chacune des étapes.

\medskip\noindent\grandscarreaux{\linewidth}{5.6cm}

\end{document}